\subsubsection{An intrinsic tame category without a model-existence hypothesis}

The cover category in Theorem~\ref{A6-tame-specialization} has an
intrinsic definition even when no proper log-regular model has been
chosen or proved to exist. This is a clarification of the category's
scope; it asserts no existence or reconstruction of a geometric model.

\begin{proposition}[Proved: an intrinsic tame Galois category]\label{A6-intrinsic-tame}
Let $R$ be a complete noetherian discrete valuation ring, let
$F=\operatorname{Frac}(R)$ have characteristic zero, and let $U/F$ be
nonempty, smooth, geometrically connected, separated and of finite type.
A test trait is a discrete valuation ring $A$, not a field, with an
$R$-algebra structure, together with an $R$-morphism
\[
 f_A:\eta_A=\operatorname{Spec}\operatorname{Frac}(A)\longrightarrow U.
\]
The ring $A$ need not be excellent, henselian, or of finite type over
$R$; the structure map need not be local. In particular, tests for
which a uniformizer of $R$ becomes a unit in $A$ are included.
Give $T_A=\operatorname{Spec}A$ the divisorial log structure of its
closed point, denoted $T_A^{\log}$.

Define $\operatorname{Tame}^{\rm tr}_R(U)$ to be the full subcategory
of $\operatorname{F\acute Et}(U)$ consisting of covers $V\to U$ such
that, for every test trait, $f_A^*V\to\eta_A$ is the generic fibre of
a \emph{finite Kummer-\'etale cover} of $T_A^{\log}$.
Then the following hold.
\begin{enumerate}
\item With the restriction of any geometric fibre functor on $U$,
 $\operatorname{Tame}^{\rm tr}_R(U)$ is a Galois category. It is closed
 in $\operatorname{F\acute Et}(U)$ under finite limits, finite sums,
 subobjects and quotients. Its fundamental group
 $\Pi^{\rm tr}_R(U)$ is a quotient of $\pi_1(U)$.
\item Pullback along an $F$-morphism $U'\to U$ between such varieties
 preserves these categories. Pullback of tame covers of
 $\operatorname{Spec}R$ gives a natural surjection
 $\Pi^{\rm tr}_R(U)\twoheadrightarrow G_F^t$.
\item Whenever a proper integral noetherian fs log-regular model
 $\mathcal X^{\log}\to(\operatorname{Spec}R)^{\log}$ with interior
 $U$ exists, one has equality of full subcategories
 \[
  \operatorname{Tame}^{\rm tr}_R(U)
     =\operatorname{Tame}_{\mathcal X}(U)
       \subseteq\operatorname{F\acute Et}(U).
 \]
 Consequently the equivalence with finite Kummer-\'etale covers of
 the strict special log fibre in
 Theorem~\ref{A6-tame-specialization} realizes this intrinsic category.
\end{enumerate}
\end{proposition}
\begin{proof}
We first verify the local categorical point without a numerical
ramification criterion. The log trait $T_A^{\log}$ is a connected
noetherian fs log-regular scheme: a uniformizer gives the chart
$\mathbb N\to A$. By
\cite[Appendix, Theorems~A.1--A.2, pp.~61--65]{HoshiLog2009A567},
finite Kummer-\'etale covers form a Galois category and restriction
to a geometric generic point is a fibre functor. Generic restriction
induces a continuous homomorphism
\[
 \rho_A:G_{\operatorname{Frac}(A)}
          \longrightarrow\pi_1^{\mathrm{k\acute et}}(T_A^{\log}).
\]
This homomorphism is surjective. Indeed, a nonempty connected finite
Kummer-\'etale cover is surjective by
\cite[Proposition~A.6, p.~56]{HoshiLog2009A567}, and is log regular,
hence normal, as recalled in the proof of Proposition~A.10 there.
Its underlying scheme is finite over the affine trait, connected
and normal, hence integral. Its generic fibre is therefore nonempty
and connected. The connected-cover criterion for surjectivity of
fundamental groups applies.

Put $J_A=\ker\rho_A$. Galois-category duality now identifies the
essential image of generic restriction with the finite continuous
$G_{\operatorname{Frac}(A)}$-sets on which $J_A$ acts trivially.
In particular, this essential image is closed under all finite
Galois-category operations, and generic restriction is fully faithful.
No claim about finiteness of a normalization of an arbitrary field
extension of a nonexcellent DVR is used.

For completeness, this gives an explicit global quotient. Fix a
geometric point $\overline u$ of $U$, write $\Pi=\pi_1(U,\overline u)$,
and, for each test, choose a geometric generic point and a path to
$\overline u$. The map $f_A$ then gives
\[
 \alpha_A:G_{\operatorname{Frac}(A)}\longrightarrow\Pi,
\]
well defined up to conjugacy. Let $N$ be the closed normal subgroup
of $\Pi$ generated by all $\alpha_A(J_A)$.
A finite $\Pi$-set satisfies every test precisely when $N$ acts
trivially. Thus
\[
 \operatorname{Tame}^{\rm tr}_R(U)
   \simeq\{\text{finite continuous }\Pi/N\text{-sets}\},
 \qquad \Pi^{\rm tr}_R(U)=\Pi/N.
\]
Normal closure removes the path choices. Equivalently, the condition
defines an intersection of full subcategories of the essentially
small category of finite $\Pi$-sets; it requires no set of moduli
parameters for the test traits. This proves the first assertion.

For an $F$-morphism $U'\to U$, composing every test for $U'$ with
this morphism proves preservation by pullback. Every test also has
a natural log morphism
$T_A^{\log}\to(\operatorname{Spec}R)^{\log}$: a boundary function
in $R$ pulls back to a nonzero element of $A$, possibly a unit.
Hence a finite Kummer-\'etale cover of the base passes every test.
The tame base-cover category is the category of finite $G_F^t$-sets,
as in Theorem~\ref{A6-tame-specialization}. A connected finite tame
extension of $F$ remains connected over the geometrically connected
$U$. The connected-cover criterion therefore proves the asserted
surjection to $G_F^t$.

Now suppose the proper log-regular model exists. If $V\to U$ belongs
to $\operatorname{Tame}_{\mathcal X}(U)$, Proposition~A.10 of the
same source extends it to a finite Kummer-\'etale cover of
$\mathcal X^{\log}$. For any test trait, the valuative criterion
for properness extends $\eta_A\to U\hookrightarrow\mathcal X$
uniquely to $T_A\to\mathcal X$ over $R$. The log structure of a
log-regular scheme is its divisorial log structure from its interior.
A local boundary-log section on $\mathcal X$ is invertible on $U$;
its pullback is therefore a nonzero element of $A$. This extends
the scheme morphism to a log morphism
$T_A^{\log}\to\mathcal X^{\log}$. Kummer-\'etale base change
supplies the required cover of the test trait. Thus
$\operatorname{Tame}_{\mathcal X}(U)
\subseteq\operatorname{Tame}^{\rm tr}_R(U)$.

Conversely, take $V$ passing all tests. For every codimension-one
boundary point $\xi$ of $\mathcal X$, the normal local ring
$A=\mathcal O_{\mathcal X,\xi}$ is a DVR and its generic point maps
to $U$. It is one of the tests, including when $\xi$ lies on the
horizontal boundary. Its generic cover extends Kummer-\'etale, hence
is tame at $\xi$ by Proposition~A.10 applied to this trait.
The normalization of $\mathcal X$ in $V$ is consequently tame at
every codimension-one boundary point. Here $\mathcal X$ is excellent,
since it is of finite type over the complete noetherian local ring
$R$, so normalization finiteness causes no issue. Proposition~A.10
on $\mathcal X^{\log}$ gives the reverse inclusion. The special-fibre
realization is now the already proved comparison theorem.
\end{proof}

The uniformizer-framed pullback extension of
Theorem~\ref{A6-tame-specialization} may therefore be formed using
$\Pi^{\rm tr}_R(U)$ for every such $U$, before a model is available.
This supplies an intrinsic arithmetic cover category and group.
A geometric special-log-fibre interpretation still requires an
existing model, and no reconstruction of $U$ from this tame quotient
is asserted. The definition deliberately uses actual finite
Kummer-\'etale extension, rather than imposing only numerical
ramification-index and residue-degree conditions on arbitrary traits.
